\chapter{量子力学的基本原理}

\textbf{2-1}\quad 证明空间反演算符 $\hat{P}\,[\hat{P}\psi(x)=\psi(-x)]$ 是厄米算符. 指出在什么条件下, $\hat{p}=-\mathrm{i}\hbar\frac{\mathrm{d}}{\mathrm{d}x}$ 是厄米算符.

\vfill

\newpage

\textbf{2-2}\quad 动量在径向方向的分量定义为 $\hat{p}_{r} = \frac{1}{2}\left(\hat{\pmb{p}}\cdot \frac{\pmb{r}}{r} + \frac{\pmb{r}}{r}\cdot \hat{\pmb{p}}\right)$ . 求出 $\hat{p}_{r}$ 在球坐标中的表达式.

\vfill

\newpage

\textbf{2-3}\quad 证明 $[\hat{p}_{x}, f(x)] = -\mathrm{i}\hbar \frac{\partial}{\partial x} f(x), [x, f(\hat{p}_{x})] = \mathrm{i}\hbar \frac{\partial}{\partial \hat{p}_{x}} f(\hat{p}_{x})$

\vfill

\newpage

\textbf{2-4}\quad 设算符 $\hat{A}$ 满足条件 $\hat{A}^{2} = 1$ ，证明 $\mathrm{e}^{\mathrm{i}\alpha \hat{A}} = \cos \alpha + \mathrm{i}\sin \alpha \hat{A}$ ，其中 $\alpha$ 为实常数。

\vfill

\newpage

\textbf{2-5}\quad 设算符 $\hat{K} = \hat{L}\hat{M}$ , $\hat{L}\hat{M} - \hat{M}\hat{L} = 1$ , 又设 $\varphi$ 为 $\hat{K}$ 的本征矢, 相应本征值为 $\lambda$ . 证明 $u \equiv \hat{L}\varphi$ 和 $v \equiv \hat{M}\varphi$ 也是 $\hat{K}$ 的本征矢, 并求出相应的本征值.

\vfill

\newpage

\textbf{2-6}\quad 粒子作一维运动， $\hat{H} = \frac{\hat{p}^{2}}{2\mu} + V(x)$ ，定态波函数为 $|n\rangle: \hat{H}|n\rangle = E_{n} |n\rangle, n = 1, 2, \cdots$

(1) 证明

\begin{equation}\tag{1}
\langle n | \hat{p} | m \rangle = a_{nm} \langle n | x | m \rangle
\end{equation}

并求系数 $a_{nm}$ .

(2) 利用式(1)推导求和公式

\begin{equation}\tag{2}
\sum_{n} (E_{n} - E_{m})^{2} \left|\langle n | x | m \rangle\right|^{2} = \frac{\hbar^{2}}{\mu^{2}} \langle m | p^{2} | m \rangle
\end{equation}

(3) 证明

\begin{equation}\tag{3}
\sum_{n} (E_{n} - E_{m}) \left|\langle n | x | m \rangle\right|^{2} = \frac{\hbar^{2}}{2\mu}
\end{equation}

\vfill

\newpage
